%% ============================================================
%% VALIDATION RESULTS -- v2, compressed to Applied Energy style.
%% ============================================================

Table~\ref{tab:validation} reports the validation against the ex-ante gates, each metric
scored against the formulation to which it applies. The baseline prescribes supply
temperature from the heating curve and does not propagate it along pipes, so downstream
temperature metrics are structurally outside its scope and are marked as such rather than as
failures.

\paragraph{Validated quantities.}
Three quantities test the formulation that produces every cost result, and all three pass.
Annual delivered energy is reproduced to \textbf{1.23}\,\% against a $2$\,\% gate; since the
network loss is the difference between generated and delivered heat, this fixes the annual
loss, which the decomposition prices. The source return temperature matches to
\textbf{2.08}\,°C (bias $-0.31$\,°C) against a $2.5$\,°C gate, measured at the one sensor not
downstream of a mixing valve. The heat-pump coefficient of performance sits in band at a mean
of \textbf{2.99}.

\paragraph{Temperature field: limitation.}
The supply-temperature field belongs to the propagating formulation and is not quantitatively
validatable on this network. The baseline does not propagate supply temperature, so its trunk
drop is identically zero against a measured $8.2$\,°C and the far-end and corridor gates do not
apply to it. On the propagating formulation the far-end error is $9.21$\,°C and is almost
entirely bias (mean absolute error and bias coincide to four decimals), matching the
$5$--$15$\,°C offset of a three-way mixing valve rather than model error; the corridor
comparison shows the same signature (bias $-7.73$\,°C). A model resolved more finely than its
metering cannot have the extra resolution verified, so the conclusions are anchored to the
annual network loss, which both formulations agree on and the measurements resolve. The far-end
node over the full annual record is reported rather than a mean over selected nodes, as the
more conservative choice; the all-node mean of $1.68$\,°C lies between the calibration- and
validation-node group means (Appendix~\ref{app:per_node_mae}).

\paragraph{Flow and energy balance.}
The source flow shows a mean absolute percentage error of \textbf{33.8}\,\%, governed by the
denominator rather than the model: the absolute error is near-constant at $12$--$13$\,m$^3$/h,
so the same error is $46$\,\% of mean flow below a quarter of peak load but $13$\,\% at half to
three-quarters load. Low-load hours are $60$\,\% of the year but carry little energy, so the
large relative error is compatible with the $1.23$\,\% annual energy match. First differences
separate a level error from a variation error, since a fixed offset cancels: the model tracks
flow level at $r=\textbf{0.91}$ and day-to-day change at $\textbf{0.80}$, and demand at $0.93$
and $0.89$. The per-step energy balance averages $6.4$\,\%, dominated by the same low-flow
hours, while the annual closure of $1.23$\,\% is well inside its gate.

\paragraph{Post-upgrade assets.}
The heat pump, electrode boiler and storage were installed after the measurement record, so
their dispatch is verified against necessary conditions rather than measurements: coefficient
of performance in band, storage state of charge within limits, and no simultaneous charge and
discharge. Two structural checks close the section: the nonlinear reference with its quadratic
constraints deactivated reproduces the baseline objective to better than $0.01$\,\%, and the
relaxation property $\mathrm{cost}(\CP) \le \mathrm{cost}(\Lone)$ holds on every one of the 135
synthetic instances.
